\Chapter{Conclusion}
\label{conclusion}

With all the combinations of extensions proposed, I would be lying if I said I understand the behavior of them all. For example, I have never fit SNET for a poisson outcome and constructed intervals using the PIPE-P method. 

Additionally, as a result of the implementation of the proposed methods in this thesis, there are many unexplored methods that are now available to use. Specifically, the biased intervals have debiased counterparts where the normal likelihoods are not re-penalized. While the crux of this thesis was arguing for biased intervals, I am not naive in thinking they will become the norm, so their debiased alternatives should be considered for their performances relative to other debiased methods. The idea of debiased PIPE CIs is not a new idea, but one that was first proposed by Biyue Dai as part of potential future work stemming from her dissertation. However, I have carried out implementation of these intervals along with the methods I have proposed. That said, her work on false discovery rate control with PIPE potentially now has the Local Quadratic Approximation alternative to consider. This is particularly of interest as she found that PIPE was generally conservative and LQA partially remedies this issue.

Another idea worth exploring is constructing the Laplace-Normal distributions under the null, $\beta_j = 0$, and then computing the p-value and understanding what properties such a quantity might have. 

My satisfaction would be a future where biased intervals are more common place and are the default for reporting in practice. While they do not provide frequentist coverage properties, they still do a good job at communicating the uncertainty in the estimates and provide insight that is undiscernible from the estimates alone. With that said, further work is needed to continue to define what properties biased intervals should have and how to evaluate them in a wider scope.

% It shouldn't be ignored that for gaussian outcomes that $\sigma^2$ is replaced by $\hat{\sigma}^2$ as a plug in estimator without making any adjustments to the distributions of $\beta$. While not a rigorous justification, there is a reasons why this was not a primary concern. Specifically, we found that the Reid estimator tends to over estimate the error variance and as such the intervals are generally over conservative for average coverage. Nevertheless, it is something that should not be excluded from future exploration.

In this dissertation I argued and proposed methods for intervals that remain consistent with the assumptions that were made when selecting a sparse penalty. Not only this, but I showed preliminary evidence that biased intervals may provide better coverage behaviors for variables selected. Additionally, I proposed three viable methods as well as how these methods can be extended to other penalties (SCAD, MCP, -NET) and other outcome types (poisson, bernoulli). With these extensions in place, I furthered the argument for biased intervals using the bias reducing penalties SCAD and MCP and constructing intervals around their estimates using the methods proposed which even further provided coverage benefits over de-biased lasso based approaches. All the methods described in this thesis have been implemented in the \texttt{pipe\_ncvreg} function in the \texttt{ncvreg} package. To my knowledge, this is the widest coverage of confidence interval construction for penalized regression approaches.

It is my hope that if the lasso were a being that it would be reading this thesis with a smile.